% Importado desde: 2026-03-28_Exploracion_Estructural_del_Protoespacio_Discreto.tex
\chapter{Exploracion estructural del protoespacio discreto: --- una formulacion inicial sin depender de Brillouin}
\label{ch:exploracion-estructural-protoespacio}






% verified-by: validators/protoespacio_minimo.py::test_linear_chain_partial_order
% verified-by: validators/protoespacio_minimo.py::test_two_cycle_unsat
% verified-by: validators/protoespacio_minimo.py::test_chain_extension
% verified-by: validators/protoespacio_minimo.py::test_cyclic_unsat

\section{Objetivo}
Queremos empezar una rama nueva del proyecto: explorar el \enquote{protoespacio discreto} sin tomar como punto de partida una zona de Brillouin ni una descripcion reciproca tipo Bloch.

La motivacion es clara:
\begin{displayquote}
\textbf{si el protoespacio es algo mas fundamental que una red cristalina efectiva, entonces su descripcion primaria no deberia depender necesariamente de un lenguaje espectral.}
\end{displayquote}

Este documento no intenta dar todavia el modelo final. Su meta es mas humilde y mas importante:
\begin{enumerate}[label=\arabic*.]
    \item proponer una descripcion estructural minima del protoespacio;
    \item distinguir datos primarios de herramientas derivadas;
    \item y dejar un marco donde el analisis de Brillouin aparezca, si aparece, solo como representacion secundaria.
\end{enumerate}

\section{Punto de partida metodologico}
Hasta ahora trabajamos con modelos periodicos porque eran la puerta mas natural para construir fermiones tipo Dirac y Weyl. Eso estuvo bien. Pero ahora conviene separar dos planos:
\begin{enumerate}[label=\arabic*.]
    \item el \textbf{plano instrumental}: redes periodicas, Hamiltonianos de Bloch, zonas de Brillouin;
    \item el \textbf{plano ontologico-estructural}: que objeto discreto podria sostener una dinamica relativista emergente.
\end{enumerate}

\begin{displayquote}
\textbf{Hipotesis de trabajo.} El lenguaje de Brillouin puede ser una muy buena herramienta para algunos modelos, pero no tiene por que ser la forma primaria de describir el protoespacio.
\end{displayquote}

\section{Que deberia ser primario}
Si dejamos de pensar primero en el espacio reciproco, los candidatos naturales a datos primarios son otros.

\subsection{Conjunto de sitios o eventos}
Debe existir un conjunto discreto base:
\begin{equation}
X=\{x_a\}.
\end{equation}

Segun la interpretacion, \(X\) puede leerse como:
\begin{enumerate}[label=\arabic*.]
    \item sitios de una red;
    \item vertices de un grafo;
    \item eventos de una estructura causal discreta;
    \item o celdas de una discretizacion mas abstracta.
\end{enumerate}

\subsection{Relacion de vecindad o incidencia}
Ademas de los puntos, hace falta una regla que diga que elementos interactuan localmente:
\begin{equation}
\cN(x_a)\subset X.
\end{equation}

Esto es mas primario que el momento de Bloch. Antes de hablar de Fourier, hay que saber quien es vecino de quien.

\subsection{Espacio interno local}
En cada sitio debe vivir un espacio interno:
\begin{equation}
\cH_{x_a}\cong \mathbb{C}^N.
\end{equation}

Ese \(N\) es lo que luego puede alojar:
\begin{enumerate}[label=\arabic*.]
    \item pseudospin;
    \item quiralidad;
    \item bloques tipo Weyl;
    \item o grados de libertad de una representacion espinorial emergente.
\end{enumerate}

\subsection{Ley dinamica local}
Hace falta tambien un operador elemental:
\begin{equation}
\cU:\bigotimes_{x_a}\cH_{x_a}\longrightarrow \bigotimes_{x_a}\cH_{x_a},
\end{equation}
o bien un Hamiltoniano local si el tiempo sigue siendo continuo.

Lo importante es que:
\begin{enumerate}[label=\arabic*.]
    \item actue localmente respecto de \(\cN\);
    \item preserve una nocion de microcausalidad;
    \item y sea capaz de producir un subsector infrarrojo tipo Dirac.
\end{enumerate}

\section{Propuesta minima de protoespacio}
Con esto, una definicion provisional y bastante general seria:
\begin{equation}
\cP_{\mathrm{disc}}=(X,\cN,\cH_{\mathrm{loc}},\cU,\cG_{\mathrm{micro}}),
\label{c30:eq:proto}
\end{equation}
donde:
\begin{enumerate}[label=\arabic*.]
    \item \(X\) es el conjunto discreto de puntos o eventos;
    \item \(\cN\) es la estructura de vecindad;
    \item \(\cH_{\mathrm{loc}}\) es la fibra local;
    \item \(\cU\) es la ley dinamica elemental;
    \item \(\cG_{\mathrm{micro}}\) es el grupo de simetrias exactas del objeto.
\end{enumerate}

\begin{figure}[h!]
\centering
\begin{tikzpicture}[
box/.style={draw, rounded corners, thick, align=center, minimum width=2.8cm, minimum height=0.95cm},
arr/.style={-{Latex[length=3mm]}, thick}
]
\node[box, fill=blue!8] (a) at (0,0) {\(X\)\\sitios/eventos};
\node[box, fill=green!8] (b) at (3.8,0) {\(\cN\)\\vecindad};
\node[box, fill=yellow!12] (c) at (7.6,0) {\(\cH_{\mathrm{loc}}\)\\fibra interna};
\node[box, fill=orange!10] (d) at (11.4,0) {\(\cU\)\\dinamica local};
\node[box, fill=red!8] (e) at (15.2,0) {\(\cG_{\mathrm{micro}}\)\\simetrias exactas};
\draw[arr] (a) -- (b);
\draw[arr] (b) -- (c);
\draw[arr] (c) -- (d);
\draw[arr] (d) -- (e);
\end{tikzpicture}
\caption{Datos estructurales minimos del protoespacio discreto antes de introducir herramientas espectrales.}
\end{figure}

\section{Que pasa con Brillouin en esta formulacion}
\subsection{Brillouin deja de ser primario}
En esta nueva lectura, el espacio reciproco ya no entra en la definicion del protoespacio. Solo puede aparecer \emph{despues}, si se satisfacen condiciones adicionales:
\begin{enumerate}[label=\arabic*.]
    \item homogeneidad;
    \item periodicidad;
    \item simetria traslacional suficiente;
    \item posibilidad de diagonalizar por modos tipo Bloch.
\end{enumerate}

\subsection{Conclusion importante}
Entonces:
\begin{displayquote}
\textbf{Brillouin no describe el protoespacio mismo; describe una representacion espectral de ciertos protoespacios discretos altamente regulares.}
\end{displayquote}

Esto es justo la separacion que nos hacia falta.

\section{Tres capas de descripcion}
Conviene ordenar todo en tres capas.

\subsection{Capa estructural}
Aqui solo viven
\[
(X,\cN,\cH_{\mathrm{loc}},\cU,\cG_{\mathrm{micro}}).
\]
Es la capa mas cercana al protoespacio.

\subsection{Capa espectral}
Si el sistema es suficientemente regular, se puede pasar a una descripcion espectral:
\[
\text{modos, autovalores, cuasienergia, Brillouin, bandas.}
\]

\subsection{Capa efectiva continua}
Linealizando un sector adecuado, se obtiene:
\[
\text{Dirac, Weyl, Clifford, Lorentz efectivo.}
\]

\begin{figure}[h!]
\centering
\begin{tikzpicture}[
box/.style={draw, rounded corners, thick, align=center, minimum width=4.2cm, minimum height=1.0cm},
arr/.style={-{Latex[length=3mm]}, thick}
]
\node[box, fill=blue!8] (a) at (0,0) {capa estructural\\protoespacio discreto};
\node[box, fill=green!8] (b) at (0,-2.2) {capa espectral\\bandas y modos};
\node[box, fill=yellow!12] (c) at (0,-4.4) {capa efectiva\\Dirac/Weyl continuo};
\draw[arr] (a) -- node[right] {\small solo si hay regularidad} (b);
\draw[arr] (b) -- node[right] {\small linealizacion infrarroja} (c);
\end{tikzpicture}
\caption{La descripcion reciproca aparece como capa intermedia posible, no como definicion del protoespacio.}
\end{figure}

\section{Primeras consecuencias para el proyecto}
\subsection{La red cubica ya no es toda la historia}
Si usamos una red cubica homogénea, eso solo fija una realizacion particular de \(X\) y \(\cN\). El protoespacio no debe identificarse sin mas con esa realizacion.

\subsection{La fibra interna se vuelve esencial}
La pregunta fuerte ya no es solo \enquote{que red produce un cono de Dirac?}, sino:
\begin{displayquote}
\textbf{que combinacion de vecindad, espacio interno y dinamica local produce un subsector relativista robusto?}
\end{displayquote}

\subsection{La microcausalidad pesa mas que la periodicidad}
Si lo que buscamos es un protoespacio, la propiedad mas profunda probablemente sea la localidad causal del paso, no la periodicidad del espectro.

\section{Dos modelos mentales posibles}
\subsection{Modelo A: red periodica enriquecida}
En este escenario, el protoespacio sigue siendo una red mas fibra mas dinamica local, y la descripcion de Brillouin es una representacion muy util porque la homogeneidad exacta es parte real del sistema.

\subsection{Modelo B: estructura local mas general}
En este escenario, la periodicidad solo fue una forma de arrancar. El protoespacio verdadero seria una estructura local mas general:
\begin{enumerate}[label=\arabic*.]
    \item tal vez no estrictamente periodica;
    \item tal vez definida por un grafo regular o casi regular;
    \item tal vez mas cercana a una red causal que a un cristal.
\end{enumerate}

En este segundo caso, Brillouin seria solo una aproximacion conveniente en sectores casi homogeneos.

\section{Que criterios deberia satisfacer este protoespacio}
Si dejamos de pensar primero en Brillouin, los criterios basicos pasan a ser:
\begin{enumerate}[label=\arabic*.]
    \item localidad de \(\cU\);
    \item existencia de una nocion de vecindad robusta;
    \item capacidad de sostener conos causales efectivos;
    \item espacio interno suficiente para una estructura tipo Clifford;
    \item subsector infrarrojo relativamente estable bajo deformaciones locales.
\end{enumerate}

Notese que ninguno de estos criterios exige de entrada un toro de Brillouin.

\section{Formula de trabajo revisada}
La pregunta fuerte del proyecto puede reescribirse asi:
\begin{displayquote}
\textbf{existe una estructura discreta local \((X,\cN,\cH_{\mathrm{loc}},\cU,\cG_{\mathrm{micro}})\) cuyo sector infrarrojo reproduzca Dirac en \(3+1\), y para la cual la interpretacion de protoespacio sea mas natural que en una simple red periodica?}
\end{displayquote}

Esta formulacion me parece mas fiel a lo que realmente estamos buscando ahora.

\section{Conclusion}
El cambio de perspectiva de este capitulo es importante. No invalida el trabajo previo; lo recoloca.

Lo que venimos construyendo hasta ahora puede leerse asi:
\begin{enumerate}[label=\arabic*.]
    \item empezamos con una familia de modelos periodicos porque eran tecnicamente manejables;
    \item eso nos permitio construir bloques Weyl y Dirac efectivos;
    \item ahora damos un paso atras conceptual y preguntamos por la estructura primaria que podria sostener todo eso;
    \item y en esa pregunta, Brillouin deja de ser fundamento y pasa a ser, en el mejor de los casos, una herramienta derivada.
\end{enumerate}

\begin{displayquote}
\textbf{La idea central que queda fijada es esta: el protoespacio discreto, si existe, deberia definirse primero por su localidad, su vecindad, su fibra interna y su dinamica elemental; solo despues, cuando haya suficiente regularidad, puede admitirse una lectura espectral tipo Brillouin.}
\end{displayquote}

\section{Siguiente paso natural}
Desde aqui se abren dos direcciones estructurales muy buenas:
\begin{enumerate}[label=\arabic*.]
    \item estudiar que tipo de vecindad y de simetrias locales permiten una emergencia isotropa sin imponer periodicidad exacta;
    \item o reformular la construccion discreta en lenguaje de grafo, red causal o sistema fibrado local, antes de pasar al espacio reciproco.
\end{enumerate}

La segunda de esas dos direcciones probablemente sea la mas fiel al espiritu de la pregunta por el protoespacio.
